% ============================================================
%  BrushHeroes — Documento de Pruebas (Módulo AR)
%  Pontificia Universidad Javeriana · Ingeniería de Sistemas
% ============================================================
\documentclass[12pt, a4paper]{article}

% ── Paquetes ─────────────────────────────────────────────────
\usepackage[spanish]{babel}
\usepackage[utf8]{inputenc}
\usepackage[T1]{fontenc}
\usepackage{lmodern}
\usepackage{geometry}
\usepackage{graphicx}
\usepackage{float}
\usepackage{caption}
\usepackage{subcaption}
\usepackage{xcolor}
\usepackage{titlesec}
\usepackage{fancyhdr}
\usepackage{hyperref}
\usepackage{enumitem}
\usepackage{tcolorbox}
\usepackage{booktabs}
\usepackage{array}
\usepackage{longtable}
\usepackage{parskip}
\usepackage{microtype}

\geometry{
  top=2.5cm, bottom=2.5cm,
  left=3cm,  right=2.5cm
}

% ── Paleta de colores (igual al manual) ───────────────────────
\definecolor{heroBlue}{HTML}{2563EB}
\definecolor{heroGreen}{HTML}{22C55E}
\definecolor{heroOrange}{HTML}{F97316}
\definecolor{heroRed}{HTML}{EF4444}
\definecolor{heroGray}{HTML}{64748B}
\definecolor{heroBg}{HTML}{F0F9FF}
\definecolor{heroLightGray}{HTML}{F8FAFC}

% ── Tipografía de secciones ───────────────────────────────────
\titleformat{\section}
  {\Large\bfseries\color{heroBlue}}
  {\thesection.}{0.8em}{}[\vspace{-4pt}\rule{\linewidth}{0.6pt}]
\titleformat{\subsection}
  {\large\bfseries\color{heroBlue!80!black}}
  {\thesubsection.}{0.6em}{}
\titleformat{\subsubsection}
  {\normalsize\bfseries\color{heroGray}}
  {\thesubsubsection.}{0.5em}{}
\titlespacing{\section}{0pt}{18pt}{8pt}
\titlespacing{\subsection}{0pt}{12pt}{4pt}
\titlespacing{\subsubsection}{0pt}{8pt}{2pt}

% ── Cajas de aviso ────────────────────────────────────────────
\tcbuselibrary{skins, breakable}

\newtcolorbox{notebox}[1][Nota]{
  colback=heroBg, colframe=heroBlue,
  fonttitle=\bfseries\small, title=#1,
  left=6pt, right=6pt, top=4pt, bottom=4pt,
  arc=4pt, boxrule=0.8pt
}
\newtcolorbox{tipbox}[1][Consejo]{
  colback=heroGreen!8, colframe=heroGreen!55!black,
  fonttitle=\bfseries\small, title=#1,
  left=6pt, right=6pt, top=4pt, bottom=4pt,
  arc=4pt, boxrule=0.8pt
}
\newtcolorbox{warnbox}[1][Importante]{
  colback=heroOrange!8, colframe=heroOrange,
  fonttitle=\bfseries\small, title=#1,
  left=6pt, right=6pt, top=4pt, bottom=4pt,
  arc=4pt, boxrule=0.8pt
}

% ── Encabezado y pie de página ────────────────────────────────
\pagestyle{fancy}
\fancyhf{}
\fancyhead[L]{\small\color{heroGray}BrushHeroes — Pruebas}
\fancyhead[R]{\small\color{heroGray}\leftmark}
\fancyfoot[C]{\small\color{heroGray}\thepage}
\renewcommand{\headrulewidth}{0.4pt}
\renewcommand{\headrule}{\hbox to\headwidth{\color{heroBlue!30}\leaders\hrule height \headrulewidth\hfill}}

\captionsetup{
  font=small,
  labelfont={bf,color=heroGray},
  textfont={color=heroGray},
  belowskip=4pt
}

% ── Comandos auxiliares ──────────────────────────────────────
\newcommand{\pass}{\textcolor{heroGreen}{\textbf{PASS}}}
\newcommand{\fail}{\textcolor{heroRed}{\textbf{FAIL}}}
\newcommand{\pending}{\textcolor{heroOrange}{\textbf{PEND.}}}
\newcommand{\reqid}[1]{\texttt{\small\color{heroGray}#1}}

% ── Metadatos PDF ─────────────────────────────────────────────
\hypersetup{
  colorlinks=true,
  linkcolor=heroBlue,
  urlcolor=heroBlue,
  pdftitle={BrushHeroes - Documento de Pruebas},
  pdfauthor={Equipo BrushHeroes - PUJ}
}

% ─────────────────────────────────────────────────────────────
\begin{document}

% ── PORTADA ──────────────────────────────────────────────────
\begin{titlepage}
  \centering
  \vspace*{1.5cm}

  {\small\color{heroGray} PONTIFICIA UNIVERSIDAD JAVERIANA}\\[0.2cm]
  {\small\color{heroGray} FACULTAD DE INGENIERÍA — INGENIERÍA DE SISTEMAS}\\[0.6cm]
  \rule{0.7\textwidth}{0.8pt}

  \vspace{1.8cm}

  \begin{tcolorbox}[
    colback=heroBlue!8, colframe=heroBlue,
    width=0.55\textwidth, halign=center,
    arc=8pt, boxrule=1.2pt,
    top=18pt, bottom=18pt
  ]
    \centering
    {\LARGE\bfseries\color{heroBlue} BrushHeroes}\\[4pt]
    {\small\color{heroGray} Plan y resultados de pruebas}
  \end{tcolorbox}

  \vspace{1.8cm}

  {\huge\bfseries Documento de Pruebas}\\[0.3cm]
  {\Large Módulo de Realidad Aumentada}\\[0.6cm]
  {\large\color{heroGray} Pruebas unitarias, atributos de calidad y usabilidad}

  \vspace{1.5cm}
  \rule{0.5\textwidth}{0.4pt}\\[0.8cm]

  \begin{tabular}{r l}
    \color{heroGray}Alcance:      & Minijuego de Realidad Aumentada (ARGame) \\
    \color{heroGray}Versión:      & 1.0 \\
    \color{heroGray}Fecha:        & \today \\
  \end{tabular}

  \vfill
  {\small\color{heroGray} Proyecto de Grado · Ingeniería de Sistemas\\
  Pontificia Universidad Javeriana · \the\year}
\end{titlepage}

% ── TABLA DE CONTENIDOS ───────────────────────────────────────
\renewcommand{\contentsname}{Tabla de contenidos}
\tableofcontents
\newpage

% =============================================================
\section{Introducción}
% =============================================================

\subsection{Propósito del documento}

Este documento presenta el plan de pruebas y los resultados obtenidos para el
\textbf{módulo de Realidad Aumentada (ARGame)} del videojuego educativo
BrushHeroes. Cubre tres niveles de verificación: pruebas unitarias automatizadas
sobre el código fuente, validación de los atributos de calidad establecidos en
la especificación de requerimientos, y protocolo de evaluación de usabilidad
con odontopediatras como expertos del dominio.

\subsection{Alcance}

El alcance de este documento se limita al \textbf{módulo de Realidad Aumentada}
(escena \texttt{ARGame.unity}). Los otros dos minijuegos (cepillado interactivo
e hilo dental) están a cargo de otros integrantes del equipo y se documentan
por separado.

\begin{notebox}
La cobertura de pruebas unitarias presentada se concentra en la lógica de
dominio del módulo AR: gestión de zonas bucales, sistema de puntuación,
ciclo de vida de la sesión y consultas sobre el estado de las bacterias.
La lógica que depende de hardware (cámara, ARCore, detección de rostro)
se valida mediante pruebas de integración manuales documentadas en la
sección de atributos de calidad.
\end{notebox}

\subsection{Definiciones y acrónimos}

\begin{table}[H]
  \centering
  \renewcommand{\arraystretch}{1.25}
  \begin{tabular}{>{\bfseries}p{2.7cm} p{10cm}}
    \toprule
    Término & Significado \\
    \midrule
    AR / RA       & Realidad Aumentada (\textit{Augmented Reality}) \\
    ARCore        & Plataforma de Realidad Aumentada de Google para Android \\
    QUIS          & \textit{Questionnaire for User Interaction Satisfaction}.
                    Cuestionario estandarizado para evaluar la satisfacción del
                    usuario con una interfaz de software \\
    NUnit         & \textit{Framework} de pruebas unitarias para .NET, integrado
                    en Unity Test Framework \\
    EditMode      & Modo de pruebas en Unity que ejecuta sin entrar a Play Mode;
                    apto para pruebas de lógica pura \\
    PlayMode      & Modo de pruebas que ejecuta dentro del ciclo de vida normal
                    de Unity \\
    \texttt{asmdef} & \textit{Assembly Definition} de Unity; archivo que define
                    una unidad de compilación independiente \\
    Sesión        & Una partida del minijuego AR, desde la pulsación de Iniciar
                    hasta el fin (éxito o tiempo agotado) \\
    Zona          & Una de las seis regiones de cepillado de la cavidad oral
                    modeladas por el juego \\
    \bottomrule
  \end{tabular}
\end{table}

% =============================================================
\section{Estrategia de pruebas}
% =============================================================

\subsection{Niveles de prueba}

El plan combina cuatro niveles de prueba con propósitos complementarios:

\begin{table}[H]
  \centering
  \caption{Niveles de prueba aplicados al módulo AR.}
  \renewcommand{\arraystretch}{1.3}
  \begin{tabular}{>{\bfseries}p{3.3cm} p{4.2cm} p{5cm}}
    \toprule
    Nivel & Propósito & Mecanismo \\
    \midrule
    Unitarias        & Verificar la lógica de cada componente en aislamiento &
                       Unity Test Framework + NUnit \\
    Integración      & Verificar la interacción entre componentes y con AR Foundation &
                       Sesión guiada en Editor + dispositivo \\
    Sistema          & Verificar atributos de calidad sobre el dispositivo
                       objetivo & Mediciones con \textit{profiler} y dispositivo real \\
    Aceptación       & Validar la utilidad y corrección clínica del juego &
                       Protocolo QUIS con odontopediatras \\
    \bottomrule
  \end{tabular}
\end{table}

\subsection{Tipos de prueba}

\begin{itemize}[leftmargin=1.5cm, itemsep=4pt]
  \item \textbf{Funcionales}: cubiertas principalmente por pruebas unitarias y
        de integración manual. Verifican que el comportamiento observable
        coincide con la especificación.
  \item \textbf{No funcionales (atributos de calidad)}: cubiertas por mediciones
        directas sobre dispositivo objetivo (FPS, latencia, RAM, tiempo de
        inicio) y por inspección estática del código y la arquitectura
        (mantenibilidad, seguridad).
  \item \textbf{Usabilidad}: cubiertas por el protocolo QUIS aplicado a
        odontopediatras como expertos del dominio, complementado con preguntas
        adicionales sobre la corrección pedagógica del contenido.
\end{itemize}

\subsection{Criterios de cierre}

Una versión del módulo AR se considera \textbf{lista para entrega} cuando:

\begin{enumerate}[leftmargin=1.5cm, itemsep=3pt]
  \item El 100\% de las pruebas unitarias pasan en el último \textit{commit}.
  \item Cada requerimiento de desempeño (RD-01 a RD-07) cumple su valor objetivo
        sobre el dispositivo objetivo.
  \item Los atributos de calidad (Confiabilidad, Disponibilidad, Seguridad,
        Mantenibilidad) cumplen los criterios cualitativos especificados.
  \item El puntaje global QUIS supera el umbral acordado con el tutor
        (recomendado: $\geq 6.0/9.0$ en cada subescala).
  \item No hay defectos abiertos clasificados como \textit{Críticos} o
        \textit{Bloqueantes}.
\end{enumerate}

% =============================================================
\section{Pruebas unitarias}
% =============================================================

\subsection{Marco de pruebas y arquitectura}

Las pruebas unitarias se implementan sobre \textbf{Unity Test Framework}
versión \texttt{1.1.33}, que utiliza \textbf{NUnit 3} como motor. Se ejecutan
en \textbf{EditMode} (sin entrar a Play Mode), lo que permite tiempos de
ejecución por debajo de 1 segundo para la suite completa.

Para que las pruebas accedan a los tipos del módulo AR, se introdujo la
siguiente estructura de \textit{Assembly Definitions}:

\begin{table}[H]
  \centering
  \caption{Estructura de \textit{assemblies} del módulo AR.}
  \renewcommand{\arraystretch}{1.3}
  \begin{tabular}{p{5cm} p{8cm}}
    \toprule
    \textbf{Asmdef} & \textbf{Contenido} \\
    \midrule
    \texttt{BrushHeroes.ARGame}        & Scripts de \textit{runtime} del módulo AR \\
    \texttt{BrushHeroes.ARGame.Tests}  & Pruebas unitarias (solo Editor) \\
    \bottomrule
  \end{tabular}
\end{table}

El asmdef de \textit{runtime} se configuró con \texttt{autoReferenced: true}
para que las clases de \textit{Editor} existentes en
\texttt{Assets/MiniGames/ARGame/Editor/} sigan compilando sin modificación.

\subsection{Componentes cubiertos}

\begin{table}[H]
  \centering
  \caption{Componentes del módulo AR cubiertos por pruebas unitarias.}
  \renewcommand{\arraystretch}{1.3}
  \begin{tabular}{>{\bfseries}p{4cm} p{8.5cm}}
    \toprule
    Componente & Lógica verificada \\
    \midrule
    MouthZoneManager & Mapeo de cada \texttt{MouthZone} a su \texttt{Transform}
                       en la cavidad oral \\
    ScoreManager     & Acumulación de puntaje, lógica de bonificación/penalización
                       por zona, métricas de precisión y enfoque, reseteo de sesión \\
    SessionManager   & Ciclo de vida de la sesión (inicio, avance, fin exitoso,
                       fallo), eventos públicos, configuración del límite de tiempo \\
    BacteriaSpawner  & Robustez de las consultas públicas sobre estado de zonas
                       cuando el \textit{spawner} no ha sido inicializado \\
    \bottomrule
  \end{tabular}
\end{table}

\subsection{Casos de prueba unitarios}

A continuación se enumera la totalidad de casos de prueba implementados,
organizados por componente. La columna \textit{Resultado} refleja la
ejecución más reciente de la suite.

\subsubsection{MouthZoneManager (9 casos)}

\renewcommand{\arraystretch}{1.2}
\begin{longtable}{>{\ttfamily\small}p{1cm} p{8.5cm} >{\centering\arraybackslash}p{2cm}}
  \toprule
  \textbf{ID} & \textbf{Caso de prueba} & \textbf{Resultado} \\
  \midrule
  \endhead
  MZ-01 & GetTransformForZone\_UpperLeft\_ReturnsUpperLeftTransform & \pass \\
  MZ-02 & GetTransformForZone\_UpperRight\_ReturnsUpperRightTransform & \pass \\
  MZ-03 & GetTransformForZone\_FrontUpper\_ReturnsFrontUpperTransform & \pass \\
  MZ-04 & GetTransformForZone\_LowerLeft\_ReturnsLowerLeftTransform & \pass \\
  MZ-05 & GetTransformForZone\_LowerRight\_ReturnsLowerRightTransform & \pass \\
  MZ-06 & GetTransformForZone\_FrontLower\_ReturnsFrontLowerTransform & \pass \\
  MZ-07 & GetTransformForZone\_AllSixZones\_ReturnDistinctTransforms & \pass \\
  MZ-08 & GetZoneTransform\_IsAliasOfGetTransformForZone & \pass \\
  MZ-09 & SetUp inyecta los 6 \texttt{Transform} privados vía reflexión & \pass \\
  \bottomrule
\end{longtable}

\subsubsection{ScoreManager (17 casos)}

\begin{longtable}{>{\ttfamily\small}p{1cm} p{8.5cm} >{\centering\arraybackslash}p{2cm}}
  \toprule
  \textbf{ID} & \textbf{Caso de prueba} & \textbf{Resultado} \\
  \midrule
  \endhead
  SC-01 & NewInstance\_ScoreIsZero                                    & \pass \\
  SC-02 & NewInstance\_AccuracyIsZero                                 & \pass \\
  SC-03 & NewInstance\_TargetZoneFocusRatioIsZero                     & \pass \\
  SC-04 & AddScore\_PositiveAmount\_IncrementsScore                   & \pass \\
  SC-05 & AddScore\_NegativeAmount\_DecrementsScore                   & \pass \\
  SC-06 & AddScore\_FiresOnScoreChangedEventWithNewValue              & \pass \\
  SC-07 & RegisterCleanResult\_TargetZone\_AddsTwoPoints              & \pass \\
  SC-08 & RegisterCleanResult\_WrongZoneWhileTargetHasWork\_SubtractsOnePoint & \pass \\
  SC-09 & RegisterCleanResult\_WrongZoneAfterTargetDone\_DoesNotChangeScore   & \pass \\
  SC-10 & RegisterCleanResult\_MixedSequence\_AccumulatesCorrectly    & \pass \\
  SC-11 & GetAccuracy\_NoSpawns\_ReturnsZero                          & \pass \\
  SC-12 & GetAccuracy\_AllSpawnedAreCleaned\_ReturnsOne               & \pass \\
  SC-13 & GetAccuracy\_HalfCleaned\_ReturnsHalf                       & \pass \\
  SC-14 & GetTargetZoneFocusRatio\_OnlyTargetZoneCleans\_ReturnsOne   & \pass \\
  SC-15 & GetTargetZoneFocusRatio\_HalfInTargetHalfOutside\_ReturnsHalf & \pass \\
  SC-16 & ResetSession\_ZeroesScoreAndCounters                        & \pass \\
  SC-17 & ResetSession\_FiresOnScoreChangedWithZero                   & \pass \\
  \bottomrule
\end{longtable}

\subsubsection{SessionManager (18 casos)}

\begin{longtable}{>{\ttfamily\small}p{1cm} p{8.5cm} >{\centering\arraybackslash}p{2cm}}
  \toprule
  \textbf{ID} & \textbf{Caso de prueba} & \textbf{Resultado} \\
  \midrule
  \endhead
  SM-01 & NewInstance\_IsNotRunning                                   & \pass \\
  SM-02 & NewInstance\_IsNotComplete                                  & \pass \\
  SM-03 & NewInstance\_TotalTimeLimitDefaultsTo180Seconds             & \pass \\
  SM-04 & StartSession\_SetsIsSessionRunningTrue                      & \pass \\
  SM-05 & StartSession\_FiresOnSessionStartedEvent                    & \pass \\
  SM-06 & StartSession\_FiresOnStepChangedWithFirstStep               & \pass \\
  SM-07 & StartSession\_AlreadyRunning\_DoesNotRefire                 & \pass \\
  SM-08 & StartSession\_PopulatesDefaultStepsCoveringAllSixZones      & \pass \\
  SM-09 & StartSession\_RemainingTimeStartsAtTotalTimeLimit           & \pass \\
  SM-10 & AdvanceToNextStep\_NotRunning\_NoOp                         & \pass \\
  SM-11 & AdvanceToNextStep\_IncrementsCurrentStepIndex               & \pass \\
  SM-12 & AdvanceToNextStep\_FiresOnStepChangedEvent                  & \pass \\
  SM-13 & AdvanceToNextStep\_BeyondLastStep\_EndsSessionSuccessfully  & \pass \\
  SM-14 & AdvanceToNextStep\_BeyondLastStep\_FiresOnSessionEnded      & \pass \\
  SM-15 & FailSession\_FromRunning\_StopsSessionAndMarksFailed        & \pass \\
  SM-16 & FailSession\_FromRunning\_FiresOnSessionFailedEvent         & \pass \\
  SM-17 & FailSession\_NotRunning\_NoOp                               & \pass \\
  SM-18 & FailSession\_AfterEnd\_NoOp                                 & \pass \\
  \bottomrule
\end{longtable}

\subsubsection{BacteriaSpawner (5 casos)}

\begin{longtable}{>{\ttfamily\small}p{1cm} p{8.5cm} >{\centering\arraybackslash}p{2cm}}
  \toprule
  \textbf{ID} & \textbf{Caso de prueba} & \textbf{Resultado} \\
  \midrule
  \endhead
  BS-01 & GetSpawnedCount\_UninitializedSpawner\_ReturnsZero          & \pass \\
  BS-02 & GetRemainingCount\_UninitializedSpawner\_ReturnsZero        & \pass \\
  BS-03 & GetMaxCount\_UninitializedSpawner\_ReturnsZero              & \pass \\
  BS-04 & IsZoneFullySpawned\_UninitializedSpawner\_ReturnsFalse      & \pass \\
  BS-05 & IsZoneComplete\_UninitializedSpawner\_ReturnsFalse          & \pass \\
  \bottomrule
\end{longtable}

\subsection{Resumen}

\begin{table}[H]
  \centering
  \caption{Resumen de la suite de pruebas unitarias.}
  \renewcommand{\arraystretch}{1.3}
  \begin{tabular}{>{\bfseries}p{4cm} >{\centering\arraybackslash}p{2.5cm}
                  >{\centering\arraybackslash}p{2.5cm} >{\centering\arraybackslash}p{2.5cm}}
    \toprule
    Componente & Casos & Pasan & Tasa \\
    \midrule
    MouthZoneManager  & 9  & 9  & 100\,\% \\
    ScoreManager      & 17 & 17 & 100\,\% \\
    SessionManager    & 18 & 18 & 100\,\% \\
    BacteriaSpawner   & 5  & 5  & 100\,\% \\
    \midrule
    \textbf{Total}    & \textbf{49} & \textbf{49} & \textbf{100\,\%} \\
    \bottomrule
  \end{tabular}
\end{table}

\subsection{Cómo ejecutar las pruebas}

\begin{enumerate}[leftmargin=1.5cm, itemsep=3pt]
  \item Abrir el proyecto en el Unity Editor.
  \item Menú \textbf{Window $\rightarrow$ General $\rightarrow$ Test Runner}.
  \item Seleccionar la pestaña \textbf{EditMode}.
  \item Pulsar \textbf{Run All}.
\end{enumerate}

% =============================================================
\section{Pruebas de atributos de calidad}
% =============================================================

Esta sección verifica el cumplimiento de los atributos no funcionales
especificados en la sección 3 de la \textit{Especificación de Requerimientos
de Software (SRS)} de BrushHeroes. Cada atributo se documenta con su criterio
de aceptación, el procedimiento de medición y el resultado.

\subsection{Desempeño}

Mediciones realizadas sobre el dispositivo objetivo (\textit{quad-core} 1.5\,GHz,
2\,GB de RAM, Android 8.0 o superior). Se utilizó el \textbf{Unity Profiler}
conectado por USB para captura de FPS, latencia y memoria.

\renewcommand{\arraystretch}{1.3}
\begin{longtable}{>{\ttfamily\small}p{1.5cm} p{4.5cm} p{4cm} >{\centering\arraybackslash}p{1.7cm}}
  \toprule
  \textbf{ID} & \textbf{Métrica} & \textbf{Valor objetivo} & \textbf{Resultado} \\
  \midrule
  \endhead
  RD-01 & Tasa de cuadros           & 60 FPS normal; mínimo 30 FPS garantizado & \pending \\
  RD-02 & Latencia de entrada       & Respuesta táctil $<$ 50\,ms              & \pending \\
  RD-03 & Latencia del módulo de RA & $<$ 100\,ms entre movimiento y \textit{overlay} & \pending \\
  RD-04 & Tiempo de inicio de la app & Pantalla principal $<$ 5\,s & \pending \\
  RD-05 & Inicio de sesión de RA    & ARCore activo $<$ 3\,s tras permiso & \pending \\
  RD-06 & Persistencia de datos     & Escritura a disco $<$ 1\,s & \pending \\
  RD-07 & Consumo de memoria        & RAM máxima 256\,MB en ejecución & \pending \\
  \bottomrule
\end{longtable}

\begin{notebox}
La columna \textbf{Resultado} debe completarse al ejecutar las mediciones
sobre el dispositivo objetivo. Reemplazar \pending{} por \pass{} o \fail{}
según corresponda y agregar la observación o el valor medido en una columna
adicional si se requiere mayor detalle.
\end{notebox}

\subsubsection{Procedimiento de medición}

\begin{itemize}[leftmargin=1.5cm, itemsep=4pt]
  \item \textbf{RD-01 (FPS):} con el Profiler conectado, ejecutar una sesión
        completa en Modo Dinámico mientras la cámara apunta al rostro. Promediar
        el FPS durante los 180\,s de juego. Mínimo aceptable: $\geq$\,30 FPS en el
        \textit{percentile 95}.
  \item \textbf{RD-02 (latencia táctil):} medir entre el evento de
        \texttt{Touch.phase = Began} y el primer reposicionamiento del cepillo
        virtual. Promedio sobre 30 toques.
  \item \textbf{RD-03 (latencia AR):} grabar pantalla a 60\,fps y comparar
        contra el movimiento físico de la cabeza, contando los \textit{frames}
        entre el movimiento físico y la actualización del \textit{overlay}.
  \item \textbf{RD-04 (tiempo de inicio):} cronometrar desde el toque del ícono
        de la app hasta que el menú principal acepta interacción. Promedio
        sobre 5 lanzamientos en frío.
  \item \textbf{RD-05 (inicio AR):} cronometrar desde la concesión del permiso
        de cámara hasta que ARCore reporta \texttt{TrackingState.Tracking} en al
        menos un \texttt{ARFace}.
  \item \textbf{RD-06 (persistencia):} cronometrar la operación de guardado
        en \texttt{PlayerDataManager} al regresar al menú principal.
  \item \textbf{RD-07 (memoria):} consumo máximo de RAM reportado por el
        Profiler durante una sesión completa.
\end{itemize}

\subsection{Confiabilidad}

\textbf{Criterio:} tasa de fallos $<$ 1\,\% de las sesiones iniciadas, donde se
considera fallo cualquier cierre inesperado o pérdida irrecuperable de progreso.

\begin{table}[H]
  \centering
  \renewcommand{\arraystretch}{1.3}
  \begin{tabular}{>{\bfseries}p{4cm} p{8cm}}
    \toprule
    Verificación & Procedimiento \\
    \midrule
    Auto-guardado & Iniciar una sesión, completar al menos una zona, forzar
                    el cierre de la app desde Android. Reabrir y verificar
                    que la racha y las estrellas reflejan la sesión previa. \\
    Degradación de AR & Cubrir la cámara durante una sesión activa.
                        Confirmar que la app muestra el \textit{overlay}
                        de "Buscando rostro" sin congelarse ni cerrarse. \\
    Pérdida de permiso & Revocar el permiso de cámara desde Android mientras
                         la app está abierta. Confirmar mensaje informativo
                         en lugar de cierre forzado. \\
    Estabilidad larga & Mantener una sesión en Modo Solo Guía por 15\,min
                        consecutivos. No deben observarse \textit{leaks} de
                        memoria ni caídas de FPS sostenidas. \\
    \bottomrule
  \end{tabular}
\end{table}

\subsection{Disponibilidad}

\textbf{Criterio:} 100\,\% en condiciones normales (aplicación local, sin
infraestructura externa). El módulo de RA puede degradarse si las condiciones
ambientales no se cumplen.

\begin{table}[H]
  \centering
  \renewcommand{\arraystretch}{1.3}
  \begin{tabular}{>{\bfseries}p{4cm} p{8cm}}
    \toprule
    Condición & Comportamiento esperado \\
    \midrule
    Iluminación $<$ 100\,lux & La detección de rostro puede fallar; la app debe
                               comunicarlo al usuario y permitir reintento. \\
    Dispositivo no compatible con ARCore & La carta de Cepillado AR debe estar
                                          visible pero el botón Jugar debe mostrar
                                          un mensaje informativo y \textbf{no}
                                          cerrar la app. \\
    Cámara ocupada por otra app & Mostrar mensaje claro e ingresar a un estado
                                   recuperable. \\
    Otros minijuegos & Deben permanecer accesibles aunque el módulo de RA no
                      esté disponible. \\
    \bottomrule
  \end{tabular}
\end{table}

\subsection{Seguridad}

\textbf{Criterio:} la aplicación cumple con la Ley 1581 de 2012 (Habeas Data,
Colombia). No se recopilan, almacenan, transmiten ni procesan datos personales
identificables de los menores usuarios.

\begin{table}[H]
  \centering
  \renewcommand{\arraystretch}{1.3}
  \begin{tabular}{>{\bfseries}p{4cm} p{8cm}}
    \toprule
    Verificación & Procedimiento \\
    \midrule
    Almacenamiento privado & Inspeccionar que todos los datos persistidos
                              se ubican en \textit{app-private storage}
                              de Android, inaccesible para otras apps. \\
    Sin SDKs externos      & Revisión de \texttt{Packages/manifest.json} y
                              \texttt{ProjectSettings} confirmando ausencia
                              de SDKs de analítica, publicidad o \textit{tracking}. \\
    Sin datos personales   & Inspección estática del esquema de
                              \texttt{PlayerDataManager} verificando que no
                              se almacena nombre real, fotografía, ubicación
                              ni edad exacta. \\
    Sin tráfico de red     & Captura con \textit{packet sniffer} durante
                              30\,min de uso normal. Debe registrarse cero
                              tráfico saliente desde el proceso de la app. \\
    Permisos mínimos       & Verificar que el manifiesto Android solicita
                              únicamente permisos esenciales (cámara). \\
    \bottomrule
  \end{tabular}
\end{table}

\subsection{Mantenibilidad}

\textbf{Criterio:} arquitectura modular con dependencias mínimas y bien
definidas; documentación inline en C\#; control de versiones Git con
\textit{commits} descriptivos.

\begin{table}[H]
  \centering
  \renewcommand{\arraystretch}{1.3}
  \begin{tabular}{>{\bfseries}p{4cm} p{8cm}}
    \toprule
    Aspecto & Verificación \\
    \midrule
    Modularidad & Cada módulo principal corresponde a una escena o prefab
                  independiente. El módulo AR tiene su propia carpeta
                  \texttt{Assets/MiniGames/ARGame/} y su propio
                  \textit{Assembly Definition} (\texttt{BrushHeroes.ARGame}). \\
    Independencia de contenido & Las descripciones de zonas, instrucciones
                                  y parámetros de juego se exponen en campos
                                  serializados editables desde el Inspector
                                  sin recompilar. \\
    Documentación inline & Comentarios XML en clases y métodos públicos
                           clave (verificación por muestreo). \\
    Trazabilidad Git & Historial de \textit{commits} con mensajes descriptivos
                       que reflejan el cambio realizado. \\
    Pruebas automatizadas & La suite unitaria descrita en la sección 3
                            permite detectar regresiones en la lógica de dominio. \\
    \bottomrule
  \end{tabular}
\end{table}

% =============================================================
\section{Protocolo QUIS de evaluación con odontopediatras}
% =============================================================

\subsection{Objetivo}

Evaluar la satisfacción del usuario y la corrección clínica del minijuego AR
mediante el cuestionario estandarizado \textbf{QUIS} (\textit{Questionnaire
for User Interaction Satisfaction}) aplicado a odontopediatras como expertos
del dominio, complementado con preguntas específicas sobre el contenido
pedagógico.

\subsection{Perfil de los evaluadores}

\begin{itemize}[leftmargin=1.5cm, itemsep=3pt]
  \item Odontólogos con especialización o experiencia en odontopediatría.
  \item Mínimo: 3 evaluadores. Recomendado: 5--8 evaluadores.
  \item Sin requisitos especiales sobre experiencia previa con videojuegos
        o tecnologías de RA (de hecho, una mezcla de niveles de familiaridad
        es deseable).
\end{itemize}

\subsection{Logística de la sesión}

Cada sesión individual tiene una duración estimada de \textbf{45 minutos} y
se compone de cuatro fases:

\begin{table}[H]
  \centering
  \renewcommand{\arraystretch}{1.3}
  \begin{tabular}{>{\bfseries}p{1.4cm} p{2.6cm} p{8cm}}
    \toprule
    Fase & Duración & Actividad \\
    \midrule
    1 & 5\,min  & Bienvenida, consentimiento informado, breve descripción
                 del juego y de la sesión. \\
    2 & 5\,min  & Demostración del juego por parte del facilitador. \\
    3 & 20\,min & Tareas guiadas (ver siguiente subsección) ejecutadas por
                  el evaluador con \textit{think-aloud} libre. \\
    4 & 15\,min & Aplicación del cuestionario QUIS y de las preguntas
                  específicas pedagógicas. \\
    \bottomrule
  \end{tabular}
\end{table}

\subsection{Tareas guiadas}

Durante la fase 3, el facilitador pide al evaluador realizar las siguientes
tareas en orden, sin asistencia salvo en caso de bloqueo total:

\begin{enumerate}[leftmargin=1.5cm, itemsep=3pt]
  \item Abrir la aplicación BrushHeroes y observar el menú principal.
  \item Identificar la racha, las estrellas y las misiones diarias.
  \item Ingresar al minijuego de Cepillado AR.
  \item Permitir que la app detecte el rostro.
  \item Iniciar una sesión en Modo \textbf{Solo Guía} y observarla completa.
  \item Volver al menú principal.
  \item Iniciar una sesión en Modo \textbf{Dinámico} e intentar completarla.
  \item Pausar la sesión y reanudarla.
  \item Salir al menú al finalizar la sesión.
  \item Revisar la pestaña Calendario.
\end{enumerate}

\subsection{Cuestionario QUIS adaptado}

El instrumento utiliza la versión corta de QUIS con cinco subescalas. Cada
ítem se valora en una escala de 0 a 9, donde 0 es la calificación más
desfavorable y 9 la más favorable. Se incluye una opción \textit{N/A}.

\subsubsection{Sección A — Reacciones generales al sistema}

Califique BrushHeroes según los siguientes pares semánticos:

\begin{table}[H]
  \centering
  \renewcommand{\arraystretch}{1.4}
  \begin{tabular}{p{4.5cm} c c c c c c c c c c p{4.5cm}}
    \toprule
    \textbf{Negativo} & 0 & 1 & 2 & 3 & 4 & 5 & 6 & 7 & 8 & 9 & \textbf{Positivo} \\
    \midrule
    Terrible          & $\square$ & $\square$ & $\square$ & $\square$ & $\square$ & $\square$ & $\square$ & $\square$ & $\square$ & $\square$ & Maravilloso \\
    Frustrante        & $\square$ & $\square$ & $\square$ & $\square$ & $\square$ & $\square$ & $\square$ & $\square$ & $\square$ & $\square$ & Satisfactorio \\
    Aburrido          & $\square$ & $\square$ & $\square$ & $\square$ & $\square$ & $\square$ & $\square$ & $\square$ & $\square$ & $\square$ & Estimulante \\
    Difícil           & $\square$ & $\square$ & $\square$ & $\square$ & $\square$ & $\square$ & $\square$ & $\square$ & $\square$ & $\square$ & Fácil \\
    Inadecuado        & $\square$ & $\square$ & $\square$ & $\square$ & $\square$ & $\square$ & $\square$ & $\square$ & $\square$ & $\square$ & Adecuado \\
    Rígido            & $\square$ & $\square$ & $\square$ & $\square$ & $\square$ & $\square$ & $\square$ & $\square$ & $\square$ & $\square$ & Flexible \\
    \bottomrule
  \end{tabular}
\end{table}

\subsubsection{Sección B — Pantalla}

\begin{enumerate}[leftmargin=1.5cm, itemsep=4pt]
  \item Lectura de los textos en pantalla: \textit{difícil} (0) — \textit{fácil} (9).
  \item Resaltado de elementos importantes: \textit{ineficaz} (0) — \textit{eficaz} (9).
  \item Organización de la información: \textit{confusa} (0) — \textit{clara} (9).
  \item Secuencia entre pantallas: \textit{confusa} (0) — \textit{clara} (9).
\end{enumerate}

\subsubsection{Sección C — Terminología e información del sistema}

\begin{enumerate}[leftmargin=1.5cm, itemsep=4pt]
  \item Uso de la terminología odontológica: \textit{inconsistente} (0) — \textit{consistente} (9).
  \item Adecuación de la terminología para niños de 4--10 años:
        \textit{inadecuada} (0) — \textit{adecuada} (9).
  \item Posición de los mensajes en pantalla: \textit{inconsistente} (0) — \textit{consistente} (9).
  \item Mensajes que indican el progreso del usuario: \textit{insuficientes} (0) — \textit{suficientes} (9).
  \item Mensajes de error (cuando aplican): \textit{poco útiles} (0) — \textit{útiles} (9).
\end{enumerate}

\subsubsection{Sección D — Aprendizaje}

\begin{enumerate}[leftmargin=1.5cm, itemsep=4pt]
  \item Aprender a operar el sistema: \textit{difícil} (0) — \textit{fácil} (9).
  \item Explorar las funciones por ensayo y error: \textit{desalentador} (0) — \textit{alentador} (9).
  \item Recordar nombres y uso de comandos: \textit{difícil} (0) — \textit{fácil} (9).
  \item Realizar las tareas es directo: \textit{nunca} (0) — \textit{siempre} (9).
  \item Mensajes de ayuda en pantalla: \textit{poco útiles} (0) — \textit{útiles} (9).
\end{enumerate}

\subsubsection{Sección E — Capacidades del sistema}

\begin{enumerate}[leftmargin=1.5cm, itemsep=4pt]
  \item Velocidad del sistema: \textit{muy lenta} (0) — \textit{adecuada} (9).
  \item Confiabilidad del sistema: \textit{poco confiable} (0) — \textit{confiable} (9).
  \item Diseñado para todos los niveles de usuario: \textit{nunca} (0) — \textit{siempre} (9).
  \item Corregir errores propios: \textit{difícil} (0) — \textit{fácil} (9).
\end{enumerate}

\subsection{Preguntas específicas sobre contenido pedagógico}

Las siguientes preguntas son adicionales al cuestionario QUIS estándar y
están diseñadas para aprovechar el conocimiento clínico de los evaluadores.

\begin{enumerate}[leftmargin=1.5cm, itemsep=5pt]
  \item ¿La técnica de cepillado mostrada en el \textbf{Modo Solo Guía}
        corresponde a la técnica recomendada para niños de 4--10 años?
        \textit{Si no, indique qué ajustaría.}
  \item ¿Las seis zonas modeladas (superior izquierda/derecha/frontal,
        inferior izquierda/derecha/frontal) cubren adecuadamente la cavidad
        oral para los objetivos pedagógicos del juego?
  \item ¿La duración de 180 segundos en \textbf{Modo Dinámico} es adecuada
        respecto al tiempo de cepillado recomendado para niños?
  \item ¿La narrativa visual (bacterias que aparecen y se eliminan) es
        apropiada y motivadora para el grupo etario objetivo?
        ¿Genera ansiedad o miedo?
  \item ¿Recomendaría este juego como herramienta de refuerzo para sus
        pacientes pediátricos? Marque una opción y comente brevemente.

  \begin{itemize}[leftmargin=2cm]
    \item[$\square$] Sí, sin reservas.
    \item[$\square$] Sí, con los ajustes mencionados.
    \item[$\square$] No por las razones que se indican.
  \end{itemize}
  \item Comentarios libres y sugerencias:\\[2cm]
\end{enumerate}

\subsection{Cálculo del puntaje QUIS}

Para cada subescala (A--E) se calcula el promedio aritmético de sus ítems.
El \textbf{puntaje global} es el promedio de las cinco subescalas. El
proyecto considera \textbf{aceptable} un puntaje global $\geq$ 6.0 / 9.0
con ninguna subescala individual por debajo de 5.0.

\subsection{Plantilla de tabulación de resultados}

\begin{table}[H]
  \centering
  \caption{Plantilla para consolidar los resultados QUIS.}
  \renewcommand{\arraystretch}{1.4}
  \begin{tabular}{>{\bfseries}p{4cm} >{\centering\arraybackslash}p{1.5cm}
                  >{\centering\arraybackslash}p{1.5cm} >{\centering\arraybackslash}p{1.5cm}
                  >{\centering\arraybackslash}p{1.5cm} >{\centering\arraybackslash}p{1.5cm}}
    \toprule
    Subescala & Ev.\,1 & Ev.\,2 & Ev.\,3 & \ldots & Promedio \\
    \midrule
    A. Reacciones generales      &  &  &  &  &  \\
    B. Pantalla                  &  &  &  &  &  \\
    C. Terminología e información &  &  &  &  &  \\
    D. Aprendizaje               &  &  &  &  &  \\
    E. Capacidades del sistema   &  &  &  &  &  \\
    \midrule
    \textbf{Puntaje global}      &  &  &  &  &  \\
    \bottomrule
  \end{tabular}
\end{table}

% =============================================================
\section{Plan de ejecución y entregables}
% =============================================================

\subsection{Cronograma sugerido}

\begin{table}[H]
  \centering
  \renewcommand{\arraystretch}{1.3}
  \begin{tabular}{>{\bfseries}p{4cm} p{4cm} p{4cm}}
    \toprule
    Actividad & Duración & Entregable \\
    \midrule
    Pruebas unitarias (CI) & Continuo  & Reporte de Test Runner \\
    Mediciones de desempeño & 1 semana  & Tabla RD-01--RD-07 con valores medidos \\
    Verificación de seguridad & 2 días   & Informe de auditoría estática \\
    Sesiones QUIS           & 2 semanas & Cuestionarios diligenciados \\
    Análisis y consolidación & 1 semana  & Versión final de este documento \\
    \bottomrule
  \end{tabular}
\end{table}

\subsection{Entregables del proceso de pruebas}

\begin{itemize}[leftmargin=1.5cm, itemsep=3pt]
  \item Este documento, actualizado con los resultados.
  \item Capturas del Test Runner mostrando los 49 casos en verde.
  \item Cuestionarios QUIS diligenciados (digitales o en papel escaneado).
  \item Tabla consolidada de puntajes QUIS.
  \item Lista de defectos detectados y su clasificación.
\end{itemize}

% =============================================================
\section{Conclusiones}
% =============================================================

\textit{[Sección a completar tras la ejecución de las mediciones de desempeño
y de las sesiones QUIS.]}

\subsection{Cobertura de pruebas alcanzada}
\textit{[Resumir cobertura unitaria y áreas no cubiertas.]}

\subsection{Cumplimiento de requerimientos}
\textit{[Resumir el cumplimiento de RD-01 a RD-07 y de los atributos
3.5.1 a 3.5.4.]}

\subsection{Resultado de la evaluación con expertos}
\textit{[Resumir puntaje QUIS y observaciones cualitativas más relevantes
de las preguntas específicas pedagógicas.]}

\subsection{Defectos identificados y trabajo futuro}
\textit{[Listar defectos abiertos y recomendaciones de mejora derivadas
de las pruebas.]}

% ── FIN DEL DOCUMENTO ────────────────────────────────────────
\vfill
\begin{center}
  \small\color{heroGray}\rule{0.5\textwidth}{0.4pt}\\[6pt]
  BrushHeroes · Documento de Pruebas · Versión 1.0\\
  Pontificia Universidad Javeriana · Ingeniería de Sistemas · \the\year
\end{center}

\end{document}
